% Deckblatt für das Template zum Erzeugen von Klausuren am AIFB
% =============================================================
%
% In dieser Datei muss im Normalfall NICHTS angepasst werden!
%
%%
\thispagestyle{empty}

\newcounter{maxAufgabenNum}
\setcounterref{maxAufgabenNum}{maxAufgaben}
\newcounter{gesPunkteNum}
\setcounterref{gesPunkteNum}{Gesamtpunkte}



%\begin{minipage}[b]{0.45\textwidth}
%\vtop{\vskip0pt\hbox{\ifthenelse{\boolean{klausurSpracheDeutsch}}{\includegraphics[height=2cm]{kit-de}}{\includegraphics[height=2cm]{}}}}
%\end{minipage}
\hfill
\begin{minipage}[b]{0.7\textwidth}
\vtop{\vskip0pt\flushright\includegraphics[height=2cm]{FAU_RWFak_Q_RGB_black}}
\end{minipage}
\begin{flushright}
\ifthenelse{\boolean{klausurSpracheDeutsch}}{{\footnotesize Institut für Wirtschaftsinformatik}}{{Lehrstuhl f{\"u}r Wirtschaftsinformatik, insb. Technische Informationssysteme}}
\end{flushright}
\begin{flushright}
\vspace{0.5cm}\klausurdatum
\end{flushright}
%~\vspace{-2cm}
%\begin{figure}[htbp]
%  \begin{minipage}{0.65\textwidth}
%    \raggedright
%    \includegraphics*[height=2cm]{kit} \\
%    \vspace{9ex} 
%  \end{minipage}
%  \begin{minipage}{0.35\textwidth}
%  	\raggedright
%  	\noindent\includegraphics*[width=1.2cm,angle=270]{aifb} \\
%  	\footnotesize{
%  	  Institut für Angewandte Informatik \\
%  	  und Formale Beschreibungsverfahren
%  	}
%  	\vspace{3ex} \\
%  	Prof. Dr. Hartmut Schmeck \\
%  	Dr.-Ing. Sanaz Mostaghim \\
%	Dipl.-Wi.-Ing. Friederike Pfeiffer % Hier Übungsleiter anpassen
%	Dipl.-Inform. Micaela W\"{u}nsche  % Hier Übungsleiter anpassen
%  \end{minipage}
%\end{figure}

\enlargethispage{3em}
\ifthenelse{\boolean{klausurSpracheDeutsch}}
{


{\Large
\begin{center}
\ifthenelse{\boolean{bonusklausur}}{Bonusklausur}{Klausur} über den Stoff der Vorlesung
	
	\textbf{"`\klausurname"'}
	
	(\klausurminuten\ Minuten)
\end{center}
}

{\large
\vspace{1cm}
\begin{tabular}{llll}
Name: & \rule{5cm}{1pt} & Vorname: & \rule{5cm}{1pt}\\
&&& \\
Matr.-Nr.: & \rule{5cm}{1pt} & Table No.: & \rule{4.5cm}{1pt}
\end{tabular}
}

~\vspace{0.0cm}

Ich bestätige, dass ich die folgenden Angaben gelesen und mich von der Vollständigkeit dieses Klausurexemplars (Seiten 1-\pageref{LastPage}) überzeugt habe.

~\vspace{0.5cm}

\begin{flushright}
\begin{minipage}{7.5cm}
\rule{\textwidth}{0.5pt}
Unterschrift des o.\,g. Klausur\-teil\-nehmers bzw. der o.\,g. Klausurteilnehmerin
\end{minipage}
\end{flushright}

\textbf{Anmerkungen:}

\begin{itemize}
	\item[1.] Legen Sie bitte Ihren Studierendenausweis bereit.
	\item[2.] Bitte tragen Sie \emph{Name}, \emph{Vorname}, \emph{Matr.-Nr.} und \emph{Tisch-Nr.} deutlich lesbar ein.
	\item[3.] Die folgenden \textbf{\ref{maxAufgaben} Aufgaben} sind vollständig zu bearbeiten.
	\item[4.] Folgende Hilfsmittel sind zugelassen: \textbf{\klausurhilfsmittel.}
	\item[5.] Täuschungsversuche führen zum Ausschluss von der Klausur.
	\item[6.] Unleserliche oder mit Bleistift geschriebene Lösungen können von der Klausur bzw. Wertung ausgeschlossen werden.
	\item[7.] Die Bearbeitungszeit beträgt \klausurminuten\ Minuten.
\end{itemize}
~
\smallbreak
{\footnotesize
\textbf{Nur für den Prüfer:}
}


}
{
{\Large
\begin{center}
\ifthenelse{\boolean{bonusklausur}}{Bonus-Exam}{Final Exam, \klausursemester}
	
	\textbf{\klausurname} 
	
	(Duration: \klausurminuten \ minutes)
\end{center}
}

{\large
\vspace{1cm}
\begin{tabular}{llll}
Last Name: & \rule{4cm}{1pt} & First Name: & \rule{4.5cm}{1pt}  \\
&&& \\
Matr.-No.: & \rule{4cm}{1pt} & Table No.: & \rule{4.5cm}{1pt} \\
\end{tabular}
}

~\vspace{0.0cm}

I confirm that I have read the following notes and that I have checked the com\-plete\-ness
of this exam (pages 1-\pageref{LastPage}).

~\vspace{0.5cm}

\begin{flushright}
\begin{minipage}{8cm}
\rule{\textwidth}{0.5pt}
Signature of the above-named exam attendee
\end{minipage}
\end{flushright}

\textbf{Notes:}

\begin{itemize}
	%\item[1.] Please have your student identity card ready.
	\item[1.] Please fill in your \emph{last name}, \emph{first name}, \emph{matriculation number} and \emph{table number} legibly.
	\item[2.] No additional materials are allowed.
	\item[3.] Attempts of deception lead to a disqualification from the exam.
	\item[4.] Unreadable answers or answers written with a pencil can be disqualified from the evaluation.
	\item[5.] You have \klausurminuten\ minutes to complete the exam.
	\item[6.] You may answer in English as well as in German.
\end{itemize}
~
\smallbreak
{\footnotesize
\textbf{Only for the examiner:}
%\hspace{2.5cm} 
%\ifthenelse{\boolean{punkteZaehlen}}
%{
%\ifnum\thegesPunkteNum=\klausurminuten
%(Maximal: \ref{Gesamtpunkte} Punkte)
%\else
%\begin{color}{red}(Maximal: \ref{Gesamtpunkte} Punkte)\end{color}
%\fi
%}
}
}

{\footnotesize

\begin{tabular}{|c|c|c|c|c|c|c|c|c|c|c|c|c|c|c|c|c|}
\hline

\ifnum\themaxAufgabenNum>0
1
\else
-
\fi
& 
\ifnum\themaxAufgabenNum>1
2 
\else
-
\fi
& 
\ifnum\themaxAufgabenNum>2
3 
\else
-
\fi
& 
\ifnum\themaxAufgabenNum>3
4 
\else
-
\fi
& 
\ifnum\themaxAufgabenNum>4
5 
\else
-
\fi
& 
\ifnum\themaxAufgabenNum>5
6 
\else
-
\fi
& 
\ifnum\themaxAufgabenNum>6
7 
\else
-
\fi
& 
\ifnum\themaxAufgabenNum>7
8 
\else
-
\fi
& 
\ifnum\themaxAufgabenNum>8
9 
\else
-
\fi
& 
\ifnum\themaxAufgabenNum>9
10
\else
-
\fi
& 
\ifnum\themaxAufgabenNum>10
11 
\else
-
\fi
& 
\ifnum\themaxAufgabenNum>11
12 
\else
-
\fi
& 
\ifnum\themaxAufgabenNum>12
13 
\else
-
\fi
& 
\ifnum\themaxAufgabenNum>13
14 
\else
-
\fi
& 
\ifnum\themaxAufgabenNum>14
15 
\else
-
\fi
& 
\ifnum\themaxAufgabenNum>15
16 
\else
-
\fi
& 
\ifthenelse{\boolean{klausurSpracheDeutsch}}{gesamt}{total}\\
\hline
%\parbox{5mm}{~\vspace{7mm}} & \parbox{5mm}{~\vspace{6mm}} & \parbox{5mm}{~\vspace{6mm}} & \parbox{5mm}{~\vspace{6mm}} & \parbox{5mm}{~\vspace{6mm}} & \parbox{5mm}{~\vspace{6mm}} & \parbox{5mm}{~\vspace{6mm}} & \parbox{5mm}{~\vspace{6mm}} & \parbox{5mm}{~\vspace{6mm}} & \parbox{5mm}{~\vspace{6mm}} & \parbox{5mm}{~\vspace{6mm}} & \parbox{5mm}{~\vspace{6mm}} & \parbox{5mm}{~\vspace{6mm}} & \parbox{5mm}{~\vspace{6mm}} & \parbox{5mm}{~\vspace{6mm}} & \parbox{5mm}{~\vspace{6mm}} & \parbox{5mm}{~\vspace{6mm}} \\

\ifnum\themaxAufgabenNum>0
\ifthenelse{\boolean{punkteZaehlen}}{\tiny($\ref{Punkte1}$)}{}
\fi
& 
\ifnum\themaxAufgabenNum>1
\ifthenelse{\boolean{punkteZaehlen}}{\tiny($\ref{Punkte2}$)}{}
\fi
& 
\ifnum\themaxAufgabenNum>2
\ifthenelse{\boolean{punkteZaehlen}}{\tiny($\ref{Punkte3}$)}{}
\fi
& 
\ifnum\themaxAufgabenNum>3
\ifthenelse{\boolean{punkteZaehlen}}{\tiny($\ref{Punkte4}$)}{}
\fi
& 
\ifnum\themaxAufgabenNum>4
\ifthenelse{\boolean{punkteZaehlen}}{\tiny($\ref{Punkte5}$)}{}
\fi
& 
\ifnum\themaxAufgabenNum>5
\ifthenelse{\boolean{punkteZaehlen}}{\tiny($\ref{Punkte6}$)}{}
\fi
& 
\ifnum\themaxAufgabenNum>6
\ifthenelse{\boolean{punkteZaehlen}}{\tiny($\ref{Punkte7}$)}{}
\fi
& 
\ifnum\themaxAufgabenNum>7
\ifthenelse{\boolean{punkteZaehlen}}{\tiny($\ref{Punkte8}$)}{}
\fi
& 
\ifnum\themaxAufgabenNum>8
\ifthenelse{\boolean{punkteZaehlen}}{\tiny($\ref{Punkte9}$)}{}
\fi
& 
\ifnum\themaxAufgabenNum>9
\ifthenelse{\boolean{punkteZaehlen}}{\tiny($\ref{Punkte10}$)}{}
\fi
& 
\ifnum\themaxAufgabenNum>10
\ifthenelse{\boolean{punkteZaehlen}}{\tiny($\ref{Punkte11}$)}{}
\fi
& 
\ifnum\themaxAufgabenNum>11
\ifthenelse{\boolean{punkteZaehlen}}{\tiny($\ref{Punkte12}$)}{}
\fi
& 
\ifnum\themaxAufgabenNum>12
\ifthenelse{\boolean{punkteZaehlen}}{\tiny($\ref{Punkte13}$)}{}
\fi
& 
\ifnum\themaxAufgabenNum>13
\ifthenelse{\boolean{punkteZaehlen}}{\tiny($\ref{Punkte14}$)}{}
\fi
& 
\ifnum\themaxAufgabenNum>14
\ifthenelse{\boolean{punkteZaehlen}}{\tiny($\ref{Punkte15}$)}{}
\fi
& 
\ifnum\themaxAufgabenNum>15
\ifthenelse{\boolean{punkteZaehlen}}{\tiny($\ref{Punkte16}$)}{}
\fi
& 
\ifthenelse{\boolean{punkteZaehlen}}
{\tiny$(
\ifnum\thegesPunkteNum=\klausurminuten
\ref{Gesamtpunkte}
\else
\begin{color}{red}\ref{Gesamtpunkte}\end{color}
\fi
)$}
{}\\


\ifnum\themaxAufgabenNum>0
\parbox{5mm}{~\vspace{5mm}}
\fi
& 
\ifnum\themaxAufgabenNum>1
\parbox{5mm}{~\vspace{5mm}}
\fi
& 
\ifnum\themaxAufgabenNum>2
\parbox{5mm}{~\vspace{5mm}}
\fi
& 
\ifnum\themaxAufgabenNum>3
\parbox{5mm}{~\vspace{5mm}}
\fi
& 
\ifnum\themaxAufgabenNum>4
\parbox{5mm}{~\vspace{5mm}}
\fi
& 
\ifnum\themaxAufgabenNum>5
\parbox{5mm}{~\vspace{5mm}}
\fi
& 
\ifnum\themaxAufgabenNum>6
\parbox{5mm}{~\vspace{5mm}}
\fi
& 
\ifnum\themaxAufgabenNum>7
\parbox{5mm}{~\vspace{5mm}}
\fi
& 
\ifnum\themaxAufgabenNum>8
\parbox{5mm}{~\vspace{5mm}}
\fi
& 
\ifnum\themaxAufgabenNum>9
\parbox{5mm}{~\vspace{5mm}}
\fi
& 
\ifnum\themaxAufgabenNum>10
\parbox{5mm}{~\vspace{5mm}}
\fi
& 
\ifnum\themaxAufgabenNum>11
\parbox{5mm}{~\vspace{5mm}}
\fi
& 
\ifnum\themaxAufgabenNum>12
\parbox{5mm}{~\vspace{5mm}}
\fi
& 
\ifnum\themaxAufgabenNum>13
\parbox{5mm}{~\vspace{5mm}}
\fi
& 
\ifnum\themaxAufgabenNum>14
\parbox{5mm}{~\vspace{5mm}}
\fi
& 
\ifnum\themaxAufgabenNum>15
\parbox{5mm}{~\vspace{5mm}}
\fi
& 
\\

\hline
\end{tabular}
}
